% ============================================================
% Atmosphäre: Zeilenabstand, Überschriften im Stil von Stationsschildern,
% Warnstreifen, Titelseite, eigene Motive (Tunnel, Strahlenzeichen, Gestalt mit Schutzmaske)
% Alle Motive sind eigene Zeichnungen (TikZ), keine Spielgrafiken.
% ============================================================
\usepackage{setspace}
\usepackage{etoolbox}

\definecolor{asphalt}{RGB}{28,28,30}
\definecolor{rust}{RGB}{150,62,25}
\definecolor{concrete}{RGB}{241,238,232}
\definecolor{glow}{RGB}{255,170,60}

\AtBeginDocument{\pagecolor{concrete}\setstretch{1.25}}
\AtBeginEnvironment{tabular}{\setstretch{1.0}}
\AtBeginEnvironment{longtable}{\setstretch{1.0}}
\addtokomafont{subsection}{\color{rust}}

% ---------- Abschnittsüberschrift (Stationsschild) ----------
\newcommand{\abschnitt}[1]{\par\vspace{1mm}%
\begin{tcolorbox}[enhanced,colback=asphalt,colframe=asphalt,coltext=white,boxrule=0pt,arc=0mm,left=8mm,top=2.5mm,bottom=2.5mm,borderline west={4mm}{0pt}{warn},borderline south={1.4pt}{0pt}{warn}]
{\Large\bfseries #1}
\end{tcolorbox}\par\vspace{1.5mm}}

% ---------- Titelleiste der Gruppenkarten (überschreibt makros.tex) ----------
\renewcommand{\gruppentitel}[5]{%
\begin{tcolorbox}[enhanced,colback=asphalt,colframe=asphalt,coltext=white,boxrule=0pt,arc=1mm,left=24mm,top=3mm,bottom=3mm,
 overlay={\fill[#1!85!black] (frame.south west) rectangle ([xshift=19mm]frame.north west);
          \node[white,font=\fontsize{30}{30}\selectfont\bfseries] at ([xshift=9.5mm]frame.west) {#2};}]
{\Large\bfseries #3}\hfill #5\\[0.8mm]\textcolor{warn}{Station: #4}
\end{tcolorbox}\par\vspace{1.5mm}}

% ---------- Warnstreifen ----------
\newcommand{\warnstreifen}[1][0.3]{\par\vspace{1mm}\noindent
\begin{tikzpicture}[x=1cm,y=1cm]
 \clip (0,0) rectangle (17.5,#1);
 \fill[warn] (0,0) rectangle (17.5,#1);
 \foreach \i in {-1,0,...,40}{\fill[asphalt] ({\i*0.8},0) -- ({\i*0.8+0.4},0) -- ({\i*0.8+0.4+#1},#1) -- ({\i*0.8+#1},#1) -- cycle;}
\end{tikzpicture}\par\vspace{1mm}}

% ---------- Strahlenzeichen (Dreiblatt) ----------
\newcommand{\trefcode}[1]{%
 \fill[#1] circle (0.13);
 \foreach \a in {90,210,330}{\fill[#1] ({\a-30}:0.22) -- ({\a-30}:0.85) arc (\a-30:\a+30:0.85) -- ({\a+30}:0.22) arc (\a+30:\a-30:0.22) -- cycle;}}
\newcommand{\trefoil}[2]{\begin{tikzpicture}[scale=#2,baseline=(current bounding box.center)]\trefcode{#1}\end{tikzpicture}}
\newcommand{\warnschild}[1]{\begin{tikzpicture}[scale=#1,baseline=(current bounding box.center)]
 \fill[warn] circle (1.05); \draw[asphalt,line width=1.6pt] circle (1.05);
 \begin{scope}[scale=0.95]\trefcode{asphalt}\end{scope}
\end{tikzpicture}}

% ---------- Gestalt mit Schutzmaske (Silhouette) ----------
\newcommand{\gestaltcode}[2]{%  Farbe, Linsenfarbe
 \fill[#1,rounded corners=5pt] (-0.75,-1.3) -- (-0.75,-0.5) -- (-0.3,-0.32) -- (0.3,-0.32) -- (0.75,-0.5) -- (0.75,-1.3) -- cycle;
 \fill[#1] (0,0.15) ellipse (0.5 and 0.62);
 \fill[#1] (-0.17,-0.62) rectangle (0.17,-0.3);
 \fill[#1] (0,-0.62) circle (0.2);
 \fill[#2] (-0.2,0.2) circle (0.14);
 \fill[#2] (0.2,0.2) circle (0.14);}
\newcommand{\gestalt}[3]{\begin{tikzpicture}[scale=#3,baseline=(current bounding box.center)]\gestaltcode{#1}{#2}\end{tikzpicture}}

% ---------- Tunnelszene (Breite 17,5 cm, Höhe frei) ----------
\newcommand{\szene}[1]{%
\begin{center}
\begin{tikzpicture}[x=1cm,y=1cm]
 \pgfmathsetmacro{\H}{#1}
 \pgfmathsetmacro{\cy}{0.56*\H}
 \clip[rounded corners=2mm] (0,0) rectangle (17.5,\H);
 \fill[black!95] (0,0) rectangle (17.5,\H);
 \foreach \k in {0,...,8}{%
   \pgfmathsetmacro{\f}{1-0.105*\k}
   \pgfmathtruncatemacro{\g}{5+8*\k}
   \fill[glow!\g!black,rounded corners=3mm] ({8.75-9.6*\f},{\cy-0.58*\H*\f}) rectangle ({8.75+9.6*\f},{\cy+0.58*\H*\f});
 }
 % Rohre an der Decke
 \foreach \s in {-1,1}{%
   \draw[black!75,line width=3pt] ({8.75-\s*8.75},{0.93*\H}) -- ({8.75-\s*0.6},{\cy+0.2});
   \draw[black!55,line width=1pt] ({8.75-\s*8.75},{0.90*\H}) -- ({8.75-\s*0.7},{\cy+0.12});
   \draw[black!70,line width=2pt] ({8.75-\s*8.75},{0.80*\H}) -- ({8.75-\s*0.9},{\cy+0.05});
 }
 % Gleise und Schwellen
 \pgfmathsetmacro{\vy}{\cy-0.25}
 \draw[black!30!white,line width=1.6pt] (4.3,0) -- (8.55,\vy);
 \draw[black!30!white,line width=1.6pt] (13.2,0) -- (8.95,\vy);
 \foreach \t in {0.06,0.17,0.3,0.46,0.66}{%
   \pgfmathsetmacro{\yy}{\t*\vy}
   \pgfmathsetmacro{\xl}{4.3+(8.55-4.3)*\t}
   \pgfmathsetmacro{\xr}{13.2-(13.2-8.95)*\t}
   \draw[black!40!white,line width=1.1pt] ({\xl-0.5},\yy) -- ({\xr+0.5},\yy);
 }
 % Gestalten im Gegenlicht
 \begin{scope}[shift={(8.2,{\cy-0.55})},scale=0.42]\gestaltcode{black}{glow!70!white}\end{scope}
 \begin{scope}[shift={(9.35,{\cy-0.5})},scale=0.36]\gestaltcode{black}{glow!70!white}\end{scope}
 % Wandlampen
 \foreach \f in {0.9,0.78,0.66,0.54,0.42}{%
   \foreach \s in {-1,1}{%
     \fill[warn,opacity=0.28] ({8.75+\s*9.6*\f*0.93},{\cy+0.58*\H*\f*0.45}) circle ({0.34*\f});
     \fill[warn] ({8.75+\s*9.6*\f*0.93},{\cy+0.58*\H*\f*0.45}) circle ({0.12*\f});
   }
 }
 % Warnschild an der Wand
 \begin{scope}[shift={(2.2,{\cy+0.25*\H})},scale=0.55]
   \fill[warn] circle (1.05); \draw[black,line width=1.6pt] circle (1.05);
   \begin{scope}[scale=0.95]\trefcode{black}\end{scope}
 \end{scope}
 % Nebel am Boden
 \shade[top color=black!0,bottom color=black!85] (0,{0.30*\H}) rectangle (17.5,0);
 % Stationsschild
 \node[fill=warn,text=black,font=\bfseries\Large,inner xsep=5mm,inner ysep=2mm] at (8.75,{\H-0.8}) {ENDSTATION 7};
\end{tikzpicture}
\end{center}}

% ---------- Titelseite ----------
\newcommand{\titelseite}{%
\szene{7.2}
\begin{center}
{\color{rust}\fontsize{34}{38}\selectfont\bfseries DIE LETZTE DEMOKRATIE\par}
\vspace{2mm}
{\large Politisches Planspiel in der Metrostation \textbf{Endstation 7}\par}
\end{center}
\warnstreifen[0.35]
\begin{center}
{\Large\bfseries \ifkurz Kurzvariante \,|\, 45 Minuten\else Vollvariante \,|\, 90 Minuten\fi}\\[1mm]
Spielleitung und Schülermaterial \,|\, inspiriert von der Idee der Inseldemokratie
\end{center}
\warnstreifen[0.35]
}
